% !PRINTED 480
\weberpage{480}

If, as above, we choose the quantity $\varrho$ as a rational function of the $\alpha$ in such a way that the $\Pi(m)$ quantities arising through the permutations of the $\alpha$,
\[
\varrho, \varrho', \varrho'' \ldots
\]
are all different from one another, then
\[
G(t) = (t - \varrho)(t - \varrho')(t - \varrho'') \ldots
\]
is a function in $\Omega$. Now each of the quantities $\varrho, \varrho', \varrho'' \ldots$ is a primitive quantity of the normal field $N = \Omega(\alpha, \alpha_1 \ldots \alpha_{m-1})$, and each of them is therefore the root of a \gesperrt{Galois} resolvent of degree $\mu$. These quantities, $\mu$ at a time, are the roots of one such equation. Hence $G(t)$, which has no equal roots, must split into nothing but irreducible factors of degree $\mu$, and from this it follows further \gesperrt{that $\mu$ is a divisor of $\Pi(m)$. At the same time $\mu$ is the degree of the Galois group} of $F(x)$.

If the degree of the \gesperrt{Galois} resolvent of an equation of degree $m$ has the greatest value $\Pi(m)$, we say, with an expression originating with \gesperrt{Kronecker}, that the equation has \gesperrt{no affect}. It has an affect that is the higher, the lower the degree $\mu$ of the \gesperrt{Galois} resolvent is. The quotient $\Pi(m) : \mu$, which must always be an integer, at most equal to $\Pi(m)$ and at least equal to 1, we shall call the \gesperrt{degree of the affect}; for an equation without affect it accordingly has the value 1. If the affect has the highest possible degree $\Pi(m)$, then the roots of the equation are themselves contained in the domain of rationality $\Omega$, and the equation is thus solved.

If by the adjunction of an algebraic quantity to $\Omega$ the \gesperrt{Galois} resolvent becomes reducible, a new resolvent of lower degree arises, and the degree of the affect of the equation $F(x) = 0$ increases. In this way, therefore, one comes nearer to the solution of the equation.

The \gesperrt{Galois} conception of the problem of solving an equation $F(x) = 0$ consists in this: by the successive adjunction of algebraic quantities of the simplest possible kind, the group is to be gradually diminished, or the affect raised, until it has reached its highest degree.

\begin{quote}
\gesperrt{The general equation of degree $m$ has no affect in the field $\Omega$ which consists of the rational functions of the coefficients $a_1, a_2 \ldots a_m$.}
\end{quote}
